\begin{center}\textbf{摘\quad 要}\end{center}

药材热风干燥中，热量与水分在物料内部同时输运并受烘房环境驱动。本文以一根细长圆柱状药材为对象，忽略两端影响，沿半径方向建立\textbf{一维径向有效扩散--导热耦合模型}，统一求解题目的四个问题，给出温度场、含水率场与达标所需时间，并对结果作物理解释与可靠性检验。由于长直径比约为六比一，热量与水分主要沿半径进出，一维径向近似是合理的；模型清楚区分题给条件与建模假设，分别处理问题一的解耦特例、问题二与三的变物性耦合、以及问题四的收缩几何。

在数值方法上，采用\textbf{节点型有限体积法}保持散度形式离散，配合隐式 BDF 时间积分。针对本题的三点组合设计构成方法上的特点：用\textbf{积分界面系数}处理随含水率、温度剧变的强非线性扩散，避免调和平均在陡峭干燥前沿处的偏差；用\textbf{材料参考坐标}把随失水收缩的动边界映射到固定区间，使方程不再显含边界移动速度；用\textbf{统一的全域最大含水率阈值判据}把场预测与达标时长的计算衔接起来。这些是针对本题的建模改进，其中有限体积、BDF 与材料坐标本身为成熟方法。

求解结果显示，问题一预热末期径向温度尚未均衡、失水集中在近表面薄层；问题二给出温度在数小时内即接近热风温度、而水分均匀化明显更慢的双时间尺度演化；问题三在同一数值解上得到达标时长约 \textbf{57.4740 h}；问题四在收缩几何下得到达标时长约 \textbf{51.0920 h}。对照三种情形可知，物性变化会延长干燥、尺寸收缩会缩短干燥，二者作用方向相反。结果的正确性由\textbf{解析级数解对照}、独立通量核算与网格一致性检验支持，说明数值实现忠实于所建方程；因缺少独立实验数据，本文不声称已验证真实药材的预测精度。

可靠性方面，网格、时间容差与界面求积加密后达标时长的变化都在很小的量级内，表明解在所设数值容差内稳定。参数扰动分析表明，在所考察的扰动范围内传质系数对达标时长影响最大，据此只能说明工况敏感程度，而非对所有参数都不敏感。模型的主要局限在于采用题给有效边界与有效物性、未显式计入相变潜热；问题四达标时刻的半径仍在给定数据内、不依赖半径外推，但与问题二、三一样都依赖环境数据四小时之后的常值外推，相关影响已单独说明。

\vspace{4pt}
\noindent\textbf{关键词：}\ 热风干燥；传热传质耦合；有效扩散系数；节点型有限体积法；移动边界；达标时长

\vspace{4pt}\hrule
